\chapter{Introduction}\label{ch:intro}

\section{WRF and WRF-Fire}

The Weather Research and Forecasting model (WRF) is a community model for atmospheric research and forecasting,
maintained by the National Center for Atmospheric Research (NCAR). This book is about its Advanced Research WRF
(ARW) dynamical core \citep{skamarock2019}, which solves the compressible, non-hydrostatic Euler equations on a
staggered grid. It uses third-order Runge--Kutta time stepping, with the acoustic modes integrated in smaller steps
\citep{wicker2002,klemp2007}. Around this core sit physics schemes for microphysics, radiation, the surface, the
land and the boundary layer. A user selects one scheme per process in a namelist.

WRF-Fire couples a model of wildland fire spread to the atmosphere \citep{coen2013,mandel2011}:
\begin{itemize}
  \item The fire lives on a finer grid than the atmosphere. Its front is the zero contour of a level-set function,
    moved at the rate of spread given by Rothermel's model \citep{rothermel1972}.
  \item The fire's heat and moisture fluxes feed back into the atmosphere, which in turn changes the winds that
    drive the fire.
  \item Since WRF~4.0 the level-set advection can use a fifth-order WENO scheme with re-initialization
    \citep{munozesparza2018}.
\end{itemize}

WRF~4.8.0 adds NCAR's Community Fire Behavior Model (CFBM), a standalone rewrite of the same algorithms
(\cref{ch:cfbm}).

\section{Why GPUs, and why it is hard}

A coupled fire forecast at 100~m resolution needs a nest of hundreds of cells on a side, sixty levels, steps of a
third of a second, and a fire grid four times finer still. The case of this book (\cref{ch:case}) takes 183{,}600
steps of the inner domain to cover 17 hours. Forecasts of this kind are useful only if they run faster than the
fire. Ensembles, which tell how much a forecast can be trusted, multiply the cost.

WRF's dynamics and most of its physics are limited by memory bandwidth, not by arithmetic. A GPU offers several times
the memory bandwidth of a CPU socket, so it is an obvious candidate (\cref{ch:gpuarch}). WRF physics on GPUs has been
studied since the port of a microphysics scheme by \citet{michalakes2008}. Moving the whole model to them is hard for
three reasons.
\begin{enumerate}
  \item \textbf{Size.} WRF~4.6.0 has about 1.9 million lines of Fortran (1.5 million without the chemistry):
    \begin{itemize}
      \item the ARW dynamics, about 74{,}000;
      \item the physics directory, about 670{,}000;
      \item WRF-Fire, about 8{,}500.
    \end{itemize}
    Even the subset one case executes is a few hundred routines.
  \item \textbf{Structure.} The code was written for CPUs:
    \begin{itemize}
      \item physics works on two-dimensional slabs;
      \item routines allocate large automatic arrays on every call;
      \item tables are initialized in the middle of the time step;
      \item data lives in a deep derived type.
    \end{itemize}
    None of this is wrong on a CPU, and all of it is costly or impossible on a GPU.
  \item \textbf{Trust.} A GPU port that changes the results cannot easily be told apart from a GPU port that is
    wrong. Fire spread is sensitive: a change in the last bit of one temperature can, after a few hours, move the fire
    front by whole cells. \Cref{ch:repro} explains this.
\end{enumerate}

\section{The approach of this project}

The port described here makes four choices.
\begin{description}
  \item[Bit-for-bit.] The GPU build must give results identical, bit for bit, to a reference CPU build of the same
    source. The reference build differs from a production build only in compiler flags and in a portable library of
    elementary functions. \Cref{ch:repro} shows how this is achieved and how it is verified.
  \item[One Fortran source.] The GPU code is OpenACC directives in the existing Fortran, under a preprocessor macro
    (\code{WRF_GPU}). With the macro undefined, the source is the CPU code. CUDA Fortran is added only where a
    profile shows that a kernel needs it. \Cref{ch:gpu} explains the choice: the project traded portability to other
    vendors for the best performance on NVIDIA A100 and H100.
  \item[Small verified steps.] The port moves one routine at a time. Each step is checked by tests that compare the
    device result with the host result of the same routine, and the whole model with the CPU reference
    (\cref{ch:port}).
  \item[Measure, then optimize.] A profiler is built with the port, so every kernel is measured from its first run.
    Optimizations that keep every bit (fusion, caching, tiling, warp-level changes) are taken only when the
    measurement says they pay (\cref{ch:prof,ch:fusion,ch:memory,ch:tiling,ch:warp}).
\end{description}

\section{How this book is organized}

\begin{description}
  \item[Part I, the WRF ecosystem:]
    \begin{itemize}
      \item the whole chain from raw data to a forecast (\cref{ch:system});
      \item the forecast case that fixes which parts of WRF are used (\cref{ch:case});
      \item how initial and boundary conditions are made (\cref{ch:icbc}).
    \end{itemize}
  \item[Part II, the ARW dynamical core:] the equations, time integration, advection, boundaries and nesting,
    turbulence, and the Coriolis and curvature terms.
  \item[Part III, physics:] how physics is coupled to the dynamics, then each scheme of the case in depth (WSM6,
    RRTMG and Dudhia radiation, the revised MM5 surface layer, Noah, YSU), and an overview of the others.
  \item[Part IV, fire:] fire behavior, the numerics of WRF-Fire, and CFBM.
  \item[Part V, software:] WRF's architecture, its parallelism on CPUs, I/O and the build.
  \item[Part VI, GPUs and the port:]
    \begin{itemize}
      \item the A100 and H100 hardware;
      \item OpenACC and CUDA Fortran;
      \item floating-point reproducibility;
      \item the porting method, and the port subsystem by subsystem.
    \end{itemize}
  \item[Part VII, performance:]
    \begin{itemize}
      \item profiling and performance models;
      \item four optimization families, each first as a concept and then on a kernel of this port;
      \item column physics on GPUs;
      \item results and multi-GPU runs.
    \end{itemize}
  \item[Part VIII, evolution:] the move from WRF~4.6.0 to 4.8.0, and an outlook.
\end{description}
The appendices hold the kernel catalogue, the case namelist, the notation, the test catalogue, a build and run
guide, and a glossary.

\paragraph{Results.} The port is in progress. Every measured number in this book comes from a file of the
repository, named next to it. A number not yet measured is marked \tbr.
